\begin{proposition}[A one-sided period estimate]\label{shared-period-bessel}
Let \(H,H_i\) be finite-dimensional complex Hilbert spaces, indexed by a
finite set, with \(D=\dim H>0\). Let \(P_i:H_i\to H\), and fix real
constants \(c_i>0\). Put
\[
 T_i=c_i^{-1/2}P_i,\qquad Q=\sum_iT_iT_i^*,\qquad
 M=\operatorname{tr}Q>0.
\]
Suppose that, for some real \(\alpha\),
\begin{equation}\label{shared-period-bessel-bound}
 \sum_i c_i^{-1}\|P_i^*h\|^2\le\alpha\|h\|^2
 \qquad(h\in H).
\end{equation}
Then \(\alpha>0\), and the centered target energy satisfies
\begin{equation}\label{shared-period-bessel-energy}
 0\le
 \eta:=\frac D{M^2}
       \left\|Q-\frac MD I_H\right\|_{\mathrm{HS}}^2
 \le\lambda:=\frac{\alpha D}{M}-1.
\end{equation}
Equivalently, the hypothesis is \(Q\le\alpha I_H\), or
\[
 \left\|\sum_iT_i v_i\right\|^2
       \le\alpha\sum_i\|v_i\|^2
       \qquad(v_i\in H_i).
\]
\end{proposition}

\begin{proof}
The first equivalence follows by evaluating the quadratic form of \(Q\).
The second uses equality of the norms of the synthesis map
\(\bigoplus_iH_i\to H\) and its adjoint. Every eigenvalue \(q_a\) of
\(Q\) lies in \([0,\alpha]\), so
\[
 \operatorname{tr}(Q^2)\le\alpha M,\qquad
 \left\|Q-\frac MD I_H\right\|_{\mathrm{HS}}^2
       =\operatorname{tr}(Q^2)-\frac{M^2}{D}.
\]
This proves the energy bound. Taking traces also gives
\(\alpha D\ge M>0\).
\end{proof}

\begin{proposition}[Restriction to the final source and target]
\label{shared-period-joint-restriction}
In Proposition~\ref{shared-period-bessel}, let \(J_i:E_i\to H_i\) be
isometries, let \(F\subset H\) have dimension \(n>0\), and let \(\Pi\)
be its orthogonal projection. Define
\[
 T_i'=T_iJ_i,\quad A_i=\Pi T_i':E_i\to F,\quad
 Q'=\sum_iT_i'T_i'^*,\quad Q_0=\sum_iA_iA_i^*,
\]
\[
 L=\operatorname{tr}(Q-Q'),\quad \rho=L/M,\quad
 \delta=(D-n)/D,\quad
 \kappa=1-\rho-(1+\lambda)\delta.
\]
If \(\kappa>0\), then
\begin{align}
 M_0:=\operatorname{tr}Q_0&\ge M\kappa>0,\label{shared-period-retained-mass}\\
 \eta_0:=
 \frac n{M_0^2}\left\|Q_0-\frac{M_0}{n}I_F\right\|_{\mathrm{HS}}^2
 &\le
 \frac{\lambda+\rho}{1-\rho-(1+\lambda)\delta}.
 \label{shared-period-retained-energy}
\end{align}
Moreover, if
\[
 \Delta=\sum_i\|T_i^*T_i-I_{H_i}\|_{\mathrm{HS}}^2,\qquad
 \Delta_0=\sum_i\|A_i^*A_i-I_{E_i}\|_{\mathrm{HS}}^2,
\]
then
\begin{equation}\label{shared-period-retained-diagonal}
 \Delta_0\le\Delta+2\alpha(D-n).
\end{equation}
If \(l=\sum_i(\dim H_i-\dim E_i)\), one may use
\begin{equation}\label{shared-period-discarded-mass}
 0\le L\le l+\sqrt{l\Delta}.
\end{equation}
\end{proposition}

\begin{proof}
Orthogonal domain restriction gives \(0\le Q'\le Q\le\alpha I_H\).
Consequently
\[
 \operatorname{tr}((I-\Pi)Q')\le\alpha(D-n),\qquad
 M_0=M-L-\operatorname{tr}((I-\Pi)Q')\ge M\kappa.
\]
On \(F\) we have \(Q_0\le\alpha I_F\). The preceding proposition gives
\[
 \eta_0\le\frac{\alpha n}{M_0}-1
 \le\frac{(1+\lambda)(1-\delta)}{\kappa}-1
 =\frac{\lambda+\rho}{\kappa}.
\]
For the diagonal estimate, work on \(\bigoplus_iE_i\) and write
\[
 G=\operatorname{diag}(T_i'^*T_i'),\qquad
 B=\operatorname{diag}(T_i'^*(I-\Pi)T_i').
\]
Then \(0\le B\le G\), the final diagonal Gram operator is \(G-B\),
and domain compression gives
\(\Delta':=\|G-I\|_{\mathrm{HS}}^2\le\Delta\).
Since \(\operatorname{tr}((G-B)B)\ge0\), even without commutation,
\[
 \begin{split}
 \Delta_0
 &=\Delta'-2\operatorname{tr}(GB)+2\operatorname{tr}B+
                          \operatorname{tr}(B^2)\\
 &\le\Delta'+2\operatorname{tr}B-\operatorname{tr}(B^2)
 \le\Delta+2\alpha(D-n).
 \end{split}
\]
Here \(\operatorname{tr}B=\operatorname{tr}((I-\Pi)Q')\).
Finally \(L\) is the trace of
\(\operatorname{diag}(T_i^*T_i)\) on the discarded domain subspace.
Subtract its identity operator and apply Cauchy--Schwarz on that
\(l\)-dimensional subspace to obtain
\eqref{shared-period-discarded-mass}.
\end{proof}

These estimates apply to the actual final maps \(A_i:E_i\to F\).
Their diagonal normalization uses
\(R'=\sum_i\dim E_i\); a frame comparison requiring
\(\sqrt{\Delta_0/R'}<1\) still requires that inequality.
One may use the smaller of \eqref{shared-period-retained-diagonal}
and any other valid diagonal bound. The lattice and height hypotheses
must refer to these same maps.
If \(\lambda=\rho=0\), the target energy stays zero under every
orthogonal target restriction. The upper Bessel estimate must therefore
be calibrated to the actual mass \(M\); a coarse operator bound alone
does not supply a small frame error.

\begin{proposition}[A signed orbital criterion]\label{shared-period-block-gram}
Use the spaces, maps and positive constants of
Proposition~\ref{shared-period-bessel}. Suppose the normalized Gram
operators have absolutely convergent operator-norm expansions
\[
 K_{ij}=c_i^{-1/2}c_j^{-1/2}P_i^*P_j
       =\sum_{t\in\mathcal T_{ij}}a_{ij,t},
 \qquad \sum_t\|a_{ij,t}\|_{\mathrm{op}}<\infty.
\]
For each \(i\), suppose a specified subset
\(\mathcal I_i\subset\mathcal T_{ii}\) has sum \(I_{H_i}\).
On \(\mathcal E=\bigoplus_iH_i\), define \(N=[N_{ij}]\) by
\[
 N_{ii}=\sum_{t\notin\mathcal I_i}a_{ii,t},
 \qquad N_{ij}=\sum_{t\in\mathcal T_{ij}}a_{ij,t}\quad(i\ne j).
\]
If, for every \(v=(v_i)\in\mathcal E\),
\begin{equation}\label{shared-period-signed-orbital-bound}
 \Re\sum_{i,j}\langle v_i,N_{ij}v_j\rangle
       \le(\alpha-1)\sum_i\|v_i\|^2,
\end{equation}
then the upper Bessel estimate
\eqref{shared-period-bessel-bound} holds.
It suffices to establish the same quadratic-form bound, with the same
\(\alpha\), for finite orbital truncations converging in operator norm
to \(N\).
\end{proposition}

\begin{proof}
Let \(S:\mathcal E\to H\) be \(S(v)=\sum_iT_iv_i\).
Absolute convergence and the diagonal identity give
\[
 S^*S=[K_{ij}]=I_{\mathcal E}+N,\qquad SS^*=Q.
\]
All rearrangements are justified because the block index set is finite.
Thus \(N\) is self-adjoint, and
\eqref{shared-period-signed-orbital-bound} says
\(S^*S\le\alpha I_{\mathcal E}\).
Since \(S^*S\ge0\), its norm is at most \(\alpha\);
\(\|SS^*\|=\|S^*S\|\) gives \(Q\le\alpha I_H\).
For the truncation assertion, pass to the limit in the quadratic form
of each fixed \(v\). Taking self-adjoint parts of the truncations
does not change their real quadratic forms.
\end{proof}

The uniform bound \eqref{shared-period-signed-orbital-bound}, calibrated
so that \(\alpha D/M-1\) is sufficiently small, is an additional
arithmetic input. Bounds for the separate squared norms
\(\|K_{ij}\|_{\mathrm{HS}}^2\) do not by themselves establish this
signed operator estimate at that scale. Off-diagonal identity cells
are retained unless their vanishing has been proved separately.
